\section{Restrictions at 24 edges}

\begin{proposition}
If a rank-six counterexample has 24 edges, then every vertex has
degree at most nine.
\end{proposition}

\begin{proof}
The one-star bound gives maximum degree at most ten. Suppose
that $v$ has degree ten and delete the ten edges through it, giving
a residual hypergraph $J$ with 14 edges and cover number at least five.
Its largest degree is at most five by the two-star bound. The sums
of the three and four largest degrees in a side are at most 11 and
13 by the four- and five-star bounds. On the side of $v$, at least
five nonisolated vertices remain and their degrees are unchanged.
All other sides retain at least six vertices by noncontainment.

There are at most ten degree-one vertices of $J$. To see this, the
other vertices of each deleted edge form a five-cover of $J$. At
most one vertex in that cover can have residual degree one. If
there were two, a single vertex could cover their at most two edges
and reduce the cover to four. Every residual degree-one vertex
occurs on at least one deleted edge because its degree in the
original hypergraph was at least two. Counting these incidences gives
the claimed bound.

Let $k$ be the number of degree-five vertices, $q\in\{0,1\}$
indicate whether one is on the side of $v$, and $k'=k-q$.
No side contains two degree-five vertices, because deleting their
stars in $J$ leaves at least five edges. Let $X_4$ count the
degree-four vertices, $W$ the degree-one vertices, and $K$ the sum
of the six side sizes in $J$. For $1\le d\le5$,
\[
 \binom d2-[d=1]=2d-3+[d=4]+3[d=5].
\]
The score $S$ that counts pair intersections with multiplicity
therefore satisfies
\[
 S=168-3K+W+X_4+3k\ \leq\ 73+X_4+3k,
\]
since $K\ge35$ and $W\le10$.
The side of $v$ has at most two degree-four vertices if $q=0$
and none if $q=1$. Indeed three such vertices violate the bound on
the three largest degrees, and a degree five and a degree four
would leave fewer than six degrees for the other three vertices
on that side. Each other side has at least six vertices. A side
there with a degree-five vertex has at least one singleton by
degree count, and a degree-four vertex on that side requires at
least two more singletons. A side there without a degree-five
vertex can have a first degree-four vertex freely, while a second
requires at least two singletons. Thus
\[
 X_4\le (5-k')+2(1-q)+\left\lfloor\frac{10-k'}2\right\rfloor.
\]
The resulting upper bounds on $S$ for $k=0,\ldots,6$ are
$85,86,88,89,91,92,93$. Intersectingness requires $S\ge91$.

Every two degree-five stars in $J$ intersect, since their union
has size at most nine. No three can have all their pairwise
codegrees equal to one. Their union would contain at least twelve
edges, giving a four-cover after covering the remaining two edges
with one vertex. Put an edge between two degree-five vertices if
their codegree is at least two. Its complement is triangle-free
and has at most $\lfloor k^2/4\rfloor$ edges by Mantel's bound.
Hence these stars alone contribute at least
$\binom k2-\lfloor k^2/4\rfloor$ to the two-side collision count $T$.
Each pair of residual edges still meets in at most four vertices,
so $T\le2(S-91)$. For $k=4,5,6$ the respective lower bounds on
$T$ are $2,4,6$ and the upper bounds are $0,2,4$.
This is a contradiction in every case.
\end{proof}

Two further constraints follow from the private pairs at degree-two
vertices. In an $m$-edge counterexample with $t$ degree-two vertices,
the resulting $t$ pairs are disjoint and leave $B=m-2t$ unpaired
edges. If part $i$ contains $t_i$ of those vertices, any other vertex
in that part is in at most one edge of each pair and in none of the
$t_i$ pairs from its own part. Thus its degree is at most $t-t_i+B$.

For distinct vertices $v,w$ of degrees at least three, suppose that
$d(w)-c(v,w)=1$. Every edge containing both $v,w$ is unpaired:
an original degree-two vertex on such an edge would have residual
degree one after deleting the star of $v$, since its private mate
avoids $v$. Together with $w$ this would give two residual
singletons in a five-cover. Therefore $c(v,w)\le d(w)-2$ whenever
$d(w)-1>B$, and the corresponding inequality holds when interchanging
$v,w$. At 24 edges, applying the mate and codegree constraints to
the degree and pair-transport enumeration leaves 265,367 degree
systems prior to the joint residual-star tests.

\begin{proposition}
Deleting any active vertex star from a rank-six counterexample leaves
two edges that meet in at least two vertices.
\end{proposition}

\begin{proof}
The residual has cover number at least five. The other five vertices
of any deleted edge cover it, so its cover number is exactly five.
Noncontainment of stars implies that every active vertex other than
the deleted vertex remains active. The residual therefore has at least
35 active vertices. If it were linear, Theorem 4.1(1) of \cite{fhmw}
would embed it into the truncated projective plane of order five,
which has only 30 vertices.
\end{proof}

In particular, form a graph on the edges of $H$ by joining pairs that
meet in at least two vertices. No vertex star of $H$ is a vertex cover
of this graph. Its vertex cover number is at least three: any one or
two edges of $H$ lie in a common star. Each one-star residual must
also have at least one excess intersection incidence over its number
of pairs.

\section{Rotation symmetry}

\begin{proposition}
Suppose an automorphism of order six cycles the six parts of $H$.
If $H$ has at most 43 edges or at most ten active vertices per part,
then $\tau(H)\le5$.
\end{proposition}

\begin{proof}
Transport the labels of one part to the others by powers of the
automorphism. Its sixth power is the identity, so the automorphism
becomes the cyclic rotation of row coordinates. When the common
part size is at most ten, pad the alphabet to ten and enumerate the
1,000,000 rows. In a counterexample distinct rows must agree in at
least one and at most four coordinates. The latter bound follows
from star noncontainment, or from the five shared vertices covering
every row. After this test within each orbit, 16,660 orbits remain.
Two distinct orbits can coexist only if all their cross pairs obey the
same test. To check compatibility it suffices to fix a representative
in the first orbit and test every row in the second.

A potential counterexample contains an orbit. A common relabeling of
all six parts commutes with rotation. The equality pattern of each
rotation of a representative therefore yields a canonical candidate
root by normalization in order of first occurrence. The program
starts once at each of the resulting 17 roots with all compatible
orbits available. To find a cover of size at most five, it branches
on the six vertices of a currently uncovered row with one fewer
available choice at each step. If the chosen orbits have a five-cover,
some remaining orbit in a counterexample must contain a row missed
by that cover. Order these orbits and branch on the first one that the
counterexample uses. Each branch retains compatible orbits that meet
the cover and compatible later orbits that break it. A five-cover of
the union of all available orbits excludes the entire branch.
The integer search in \texttt{edge24/cyclic.cpp} exhausts all roots.

Now suppose the common active part size $q$ is at least eleven and
$m\le43$. The main theorem gives $m\ge24$. Each part has the same
number of degree-two vertices, at most $b=\lfloor m/12\rfloor$ because
the degree-two pairs are disjoint. Every other active vertex has
degree at least three. This requires $m\ge33-b$, which excludes
$m<31$ in the stated range. For a fixed number of twos, convexity
maximizes the pair count by assigning all excess degree to one
vertex. The resulting upper bound increases on reducing the side
size to eleven and increasing the number of twos to $b$. The total
over all six sides is at most
\[
 6\left(\binom{m-30+b}{2}+30-2b\right).
\]
Subtract this from $\binom m2$. For each constant value of $b$ the
difference is a concave quadratic in $m$. At the endpoints 31, 35,
36, and 43 it equals 291, 313, 270, and 39, respectively. So the
pair score falls short of the number of pairs throughout the range.
\end{proof}

\section{Extensions of an affine-plane core}

Give each point of the affine plane over the field of five elements
the row of its six incident lines, one per direction. These labels are old and any
other labels are called new.

\begin{proposition}
An intersecting extension containing at least twelve unchanged
affine rows, in their original parts, has a five-cover.
\end{proposition}

\begin{proof}
It suffices to keep any twelve of the specified rows as the core.
A counterexample needs at least one new label in every direction.
Each row containing new labels must meet the core using its old
lines. Two nonparallel lines cover only nine points, so at least
three labels in the row are old. Enumerate all choices of one to three
omitted directions and all old line choices in the other directions.
For each union, mark every twelve-point subset of it as admitting new
labels in the omitted directions. There are 5,169,500 subsets that
admit new labels in all six directions. The group of affine linear
maps and translations partitions these into 482 orbits.

For each representative, use one symbolic new label per direction.
Actual equalities between new labels imply symbolic equalities. A
cover using only old labels therefore applies to every realization
of the covered symbolic types. A direction with fewer than two
singleton-new types and no multi-new type is impossible. A new
label would lie in at most one edge, since all five other labels are
fixed by the sole allowed type. This excludes 17 orbits. If a
direction has exactly two singleton-new types and no multi-new type,
both types must occur. Of the remaining orbits, 124 have an old-label
five-cover after filtering for compatibility with all such mandatory
types.

For the remaining 341 orbits, start with the core and the mandatory
types. Branch on the first available type breaking the current
old-label five-cover, retaining only compatible future types. The
search rejects a branch if all remaining types have an old-label
five-cover, if a needed direction can no longer support a new label
of degree at least two, or if a chosen old label can no longer occur
twice. For the last test an old-only type contributes one possible
occurrence, and a type with a new label contributes two, which
suffices for a threshold of two. The program never limits how many
actual rows can realize a type with new labels. Every actual extension
would therefore survive some branch. The exact search in
\texttt{edge24/affine-core.js} exhausts all branches in 18,645 nodes.
\end{proof}
